\begin{frame}{Peer-Review 체크리스트 — 5항목}
  \small
  \begin{tabular}{@{}ll@{}}
    \toprule
    항목 & 확인할 것 \\
    \midrule
    \kb{환경 설명} & 상태·행동·보상이 명확히 정의됐는가 \\
    \kb{알고리즘 선택} & 환경의 특성(이산/연속 행동, 보상의 희소성 등)에
    알고리즘이 맞았는가 \\
    \kb{수식적 정당화} & 요구사항을 충족했고, 실제로 관찰한 결과와
    일치하는가 \\
    \kb{결과의 정직성} & 수렴하지 않은 부분·학습이 불안정했던 구간을
    숨기지 않고 원인을 분석했는가 \\
    \kb{탐험-활용 균형} & Week 2에서 배운 딜레마가 실제 학습 곡선에서
    어떻게 드러났는가 \\
    \bottomrule
  \end{tabular}
  \vskip0.8em
  5항목은 16.1 채점 기준의 \textbf{축 그대로} — 리뷰어와 발표팀이
  \textbf{같은 기준표}를 공유해야 ``감점''과 ``좋은 지적''이 \textbf{같은
  언어}로 이야기될 수 있다.
\end{frame}
